\documentclass[10pt,a4paper]{amsart}
\usepackage{amsmath,amssymb,amsthm,graphicx}
% Retained minimal preamble for the Stacks Project chapter sites-cohomology.tex (Cohomology on Sites).
% Verbatim excerpt of the chapter preamble: Xy-pic setup (native preamble lines 14-18) and the native
% theorem declarations with the mathematical macros (native preamble lines 154-217). Copied mechanically;
% nothing added, removed or redefined. The wrapper class and the amsmath/amssymb/amsthm/graphicx base
% packages are supplied by the assembler, so the class/externaldocument part of the native preamble is omitted.
% For commutative diagrams we use Xy-pic
\usepackage[all]{xy}

% We use 2cell for 2-commutative diagrams.
\xyoption{2cell}
% Theorem environments.
%
\theoremstyle{plain}
\newtheorem{theorem}[subsection]{Theorem}
\newtheorem{proposition}[subsection]{Proposition}
\newtheorem{lemma}[subsection]{Lemma}

\theoremstyle{definition}
\newtheorem{definition}[subsection]{Definition}
\newtheorem{example}[subsection]{Example}
\newtheorem{exercise}[subsection]{Exercise}
\newtheorem{situation}[subsection]{Situation}

\theoremstyle{remark}
\newtheorem{remark}[subsection]{Remark}
\newtheorem{remarks}[subsection]{Remarks}

\numberwithin{equation}{subsection}

% Macros
%
\def\lim{\mathop{\mathrm{lim}}\nolimits}
\def\colim{\mathop{\mathrm{colim}}\nolimits}
\def\Spec{\mathop{\mathrm{Spec}}}
\def\Hom{\mathop{\mathrm{Hom}}\nolimits}
\def\Ext{\mathop{\mathrm{Ext}}\nolimits}
\def\SheafHom{\mathop{\mathcal{H}\!\mathit{om}}\nolimits}
\def\SheafExt{\mathop{\mathcal{E}\!\mathit{xt}}\nolimits}
\def\Sch{\mathit{Sch}}
\def\Mor{\mathop{\mathrm{Mor}}\nolimits}
\def\Ob{\mathop{\mathrm{Ob}}\nolimits}
\def\Sh{\mathop{\mathit{Sh}}\nolimits}
\def\NL{\mathop{N\!L}\nolimits}
\def\CH{\mathop{\mathrm{CH}}\nolimits}
\def\proetale{{pro\text{-}\acute{e}tale}}
\def\etale{{\acute{e}tale}}
\def\QCoh{\mathit{QCoh}}
\def\Ker{\mathop{\mathrm{Ker}}}
\def\Im{\mathop{\mathrm{Im}}}
\def\Coker{\mathop{\mathrm{Coker}}}
\def\Coim{\mathop{\mathrm{Coim}}}

% Boxtimes
%
\DeclareMathSymbol{\boxtimes}{\mathbin}{AMSa}{"02}

%
% Macros for moduli stacks/spaces
%
\def\QCohstack{\mathcal{QC}\!\mathit{oh}}
\def\Cohstack{\mathcal{C}\!\mathit{oh}}
\def\Spacesstack{\mathcal{S}\!\mathit{paces}}
\def\Quotfunctor{\mathrm{Quot}}
\def\Hilbfunctor{\mathrm{Hilb}}
\def\Curvesstack{\mathcal{C}\!\mathit{urves}}
\def\Polarizedstack{\mathcal{P}\!\mathit{olarized}}
\def\Complexesstack{\mathcal{C}\!\mathit{omplexes}}
% \Pic is the operator that assigns to X its picard group, usage \Pic(X)
% \Picardstack_{X/B} denotes the Picard stack of X over B
% \Picardfunctor_{X/B} denotes the Picard functor of X over B
\def\Pic{\mathop{\mathrm{Pic}}\nolimits}
\def\Picardstack{\mathcal{P}\!\mathit{ic}}
\def\Picardfunctor{\mathrm{Pic}}
\def\Deformationcategory{\mathcal{D}\!\mathit{ef}}

% --- appended 2026-10-12 (Chief): the project own definitions of the reference/slogan/history
% display environments, copied mechanically from the shared Stacks preamble lines 8-12
% (\usepackage{verbatim} plus the three \newenvironment lines, which the chapter preamble keeps
% verbatim from the project preamble). Nothing added, removed or redefined: these environments
% carry no mathematical content and the project itself typesets them as comments. Without these
% three lines a counted statement whose own text uses them fails with Environment slogan undefined.
\usepackage{verbatim}
\newenvironment{reference}{\comment}{\endcomment}
\newenvironment{slogan}{\comment}{\endcomment}
\newenvironment{history}{\comment}{\endcomment}
\usepackage{cleveref}
\crefname{originalsection}{original section}{original sections}
\Crefname{originalsection}{Original section}{Original sections}
\makeatletter
\newlabel{section-compare}{{0EWK}{source}{}{sourcepointer-1}{}}
\newlabel{section-compare@cref}{{[originalsection][0][]0EWK}{[1][0][]source}{}{}{}}
\makeatother
\crefname{originallemma}{original lemma}{original lemmas}
\Crefname{originallemma}{Original lemma}{Original lemmas}
\makeatletter
\newlabel{lemma-higher-direct-images}{{072W}{source}{}{sourcepointer-2}{}}
\newlabel{lemma-higher-direct-images@cref}{{[originallemma][0][]072W}{[1][0][]source}{}{}{}}
\makeatother
\crefname{originallemma}{original lemma}{original lemmas}
\Crefname{originallemma}{Original lemma}{Original lemmas}
\makeatletter
\newlabel{homology-lemma-characterize-weak-serre-subcategory}{{0754}{source}{}{sourcepointer-3}{}}
\newlabel{homology-lemma-characterize-weak-serre-subcategory@cref}{{[originallemma][0][]0754}{[1][0][]source}{}{}{}}
\makeatother
\crefname{originaldefinition}{original definition}{original definitions}
\Crefname{originaldefinition}{Original definition}{Original definitions}
\makeatletter
\newlabel{homology-definition-serre-subcategory}{{02MO}{source}{}{sourcepointer-4}{}}
\newlabel{homology-definition-serre-subcategory@cref}{{[originaldefinition][0][]02MO}{[1][0][]source}{}{}{}}
\makeatother
\crefname{originallemma}{original lemma}{original lemmas}
\Crefname{originallemma}{Original lemma}{Original lemmas}
\makeatletter
\newlabel{homology-lemma-five-lemma}{{05QB}{source}{}{sourcepointer-5}{}}
\newlabel{homology-lemma-five-lemma@cref}{{[originallemma][0][]05QB}{[1][0][]source}{}{}{}}
\makeatother
\crefname{originalsection}{original section}{original sections}
\Crefname{originalsection}{Original section}{Original sections}
\makeatletter
\newlabel{derived-section-triangulated-sub}{{06UP}{source}{}{sourcepointer-6}{}}
\newlabel{derived-section-triangulated-sub@cref}{{[originalsection][0][]06UP}{[1][0][]source}{}{}{}}
\makeatother
\crefname{originallemma}{original lemma}{original lemmas}
\Crefname{originallemma}{Original lemma}{Original lemmas}
\makeatletter
\newlabel{lemma-equivalence-unbounded-two}{{0D7W}{source}{}{sourcepointer-7}{}}
\newlabel{lemma-equivalence-unbounded-two@cref}{{[originallemma][0][]0D7W}{[1][0][]source}{}{}{}}
\makeatother
\title{Comparing two topologies: derived categories of adequate modules agree}
\author{The Stacks Project Authors}
\date{stacks-project:sites-cohomology.tex:local-snapshot-20261011; GNU FDL 1.2 or later}
\begin{document}
\maketitle
\noindent\textit{Editorial note.} All delivered mathematical text is the author's own, copied from the pinned Stacks Project chapter file sites-cohomology.tex; no statement or proof step was rewritten, reordered or removed.\par\medskip
\noindent\textit{Editorial note.} The proof is the author's complete argument: the identification of the categories, the vanishing of the higher direct images, the weak Serre property by the five lemma and the verification of the criteria are all written out. The phrases matching a generic omission filter inside the counted span are pointers to earlier discussions of notation, not disclaimers of arguments.\par\medskip
\noindent\textit{Editorial note.} Every labelled reference occurring inside the counted spans is carried in assembly.external\_references as a non-load-bearing navigation pointer (same-chapter ones at their definition site, cross-chapter ones at the invocation inside the counted span), so no load-bearing mathematics is replaced by a stub. No bibliography citation occurs inside the counted spans.\par\medskip
\noindent\textit{Editorial note.} The checked chapter list for same-name collisions is the set of stacks chapters staged in this repository: chow, cohomology, divisors, etale-cohomology, flat, more-algebra, more-morphisms, morphisms, spaces-more-morphisms, varieties together with this card's chapter sites-cohomology.\par\medskip
\noindent\textit{Editorial note.} The preamble is the chapter's own retained minimal preamble (sources/stacks-sites-cohomology/retained-preamble.tex), a verbatim excerpt of the native preamble: the Xy-pic setup and the native theorem declarations with the mathematical macros.\par\medskip
\noindent\textit{Allowed standard prerequisites (editorial).}\begin{itemize}
\item The comparison of topologies for the first time. Same chapter, Section section-compare, which identifies the morphism of ringed sites and its basic properties. Supplies the morphism epsilon and the earlier statements used in the proof.
\item Higher direct images as sheafifications of local cohomology. Same chapter, Lemma lemma-higher-direct-images. Turns the vanishing assumptions into the vanishing of the positive higher direct images.
\item The comparison criterion for derived categories. Same chapter, Lemma lemma-equivalence-unbounded-two, whose conditions are checked at the end. Supplies the equivalence of triangulated categories once the five conditions are verified.
\item The five lemma and weak Serre subcategories. Homology, Lemma homology-lemma-five-lemma and Definition homology-definition-serre-subcategory. Shows the image of the subcategory is a weak Serre subcategory.
\end{itemize}\medskip
\ifdefined\hypertarget\hypertarget{sourcepointer-1}{}\fi
\noindent\textit{Source reference.} Original section 0EWK: sources/stacks-sites-cohomology/pinned-original-source.tex, source lines 6177--6177. Navigation pointer invoked at proof line 6548; the counted span contains the pointer only and the cited result is not reproduced inside it, so the pointer is registered as a non-load-bearing navigation stub in the established form for Stacks material. No bibliography citation is involved.\par\medskip
\ifdefined\hypertarget\hypertarget{sourcepointer-2}{}\fi
\noindent\textit{Source reference.} Original lemma 072W: sources/stacks-sites-cohomology/pinned-original-source.tex, source lines 747--747. Navigation pointer invoked at proof line 6599; the counted span contains the pointer only and the cited result is not reproduced inside it, so the pointer is registered as a non-load-bearing navigation stub in the established form for Stacks material. No bibliography citation is involved.\par\medskip
\ifdefined\hypertarget\hypertarget{sourcepointer-3}{}\fi
\noindent\textit{Source reference.} Original lemma 0754: sources/stacks-sites-cohomology/pinned-original-source.tex, source lines 6603--6603. Navigation pointer invoked at proof line 6603; the counted span contains the pointer only and the cited result is not reproduced inside it, so the pointer is registered as a non-load-bearing navigation stub in the established form for Stacks material. No bibliography citation is involved.\par\medskip
\ifdefined\hypertarget\hypertarget{sourcepointer-4}{}\fi
\noindent\textit{Source reference.} Original definition 02MO: sources/stacks-sites-cohomology/pinned-original-source.tex, source lines 6628--6628. Navigation pointer invoked at proof line 6628; the counted span contains the pointer only and the cited result is not reproduced inside it, so the pointer is registered as a non-load-bearing navigation stub in the established form for Stacks material. No bibliography citation is involved.\par\medskip
\ifdefined\hypertarget\hypertarget{sourcepointer-5}{}\fi
\noindent\textit{Source reference.} Original lemma 05QB: sources/stacks-sites-cohomology/pinned-original-source.tex, source lines 6647--6647. Navigation pointer invoked at proof line 6647; the counted span contains the pointer only and the cited result is not reproduced inside it, so the pointer is registered as a non-load-bearing navigation stub in the established form for Stacks material. No bibliography citation is involved.\par\medskip
\ifdefined\hypertarget\hypertarget{sourcepointer-6}{}\fi
\noindent\textit{Source reference.} Original section 06UP: sources/stacks-sites-cohomology/pinned-original-source.tex, source lines 6658--6658. Navigation pointer invoked at proof line 6658; the counted span contains the pointer only and the cited result is not reproduced inside it, so the pointer is registered as a non-load-bearing navigation stub in the established form for Stacks material. No bibliography citation is involved.\par\medskip
\ifdefined\hypertarget\hypertarget{sourcepointer-7}{}\fi
\noindent\textit{Source reference.} Original lemma 0D7W: sources/stacks-sites-cohomology/pinned-original-source.tex, source lines 6482--6482. Navigation pointer invoked at proof line 6660; the counted span contains the pointer only and the cited result is not reproduced inside it, so the pointer is registered as a non-load-bearing navigation stub in the established form for Stacks material. No bibliography citation is involved.\par\medskip
Let $\mathcal{C}$ be a category. Let
$\text{Cov}(\mathcal{C}) \supset \text{Cov}'(\mathcal{C})$
be two ways to endow $\mathcal{C}$ with the structure of a site.
Denote $\tau$ the topology corresponding to $\text{Cov}(\mathcal{C})$
and $\tau'$ the topology corresponding to $\text{Cov}'(\mathcal{C})$.
Then the identity functor on $\mathcal{C}$ defines a morphism
of sites
$$
\epsilon : \mathcal{C}_\tau \longrightarrow \mathcal{C}_{\tau'}
$$
where $\epsilon_*$ is the identity functor on underlying presheaves and
where $\epsilon^{-1}$ is the $\tau$-sheafification of a $\tau'$-sheaf
(hence clearly exact). Let $\mathcal{O}$ be a sheaf of rings for the
$\tau$-topology. Then $\mathcal{O}$ is also a sheaf for the $\tau'$-topology
and $\epsilon$ becomes a morphism of ringed sites
$$
\epsilon :
(\mathcal{C}_\tau, \mathcal{O}_\tau)
\longrightarrow
(\mathcal{C}_{\tau'}, \mathcal{O}_{\tau'})
$$
For more discussion, see Section \ref{section-compare}.

\begin{lemma}
\label{lemma-compare-topologies-derived-adequate-modules}
With $\epsilon : (\mathcal{C}_\tau, \mathcal{O}_\tau) \to
(\mathcal{C}_{\tau'}, \mathcal{O}_{\tau'})$ as above.
Let $\mathcal{B} \subset \Ob(\mathcal{C})$ be a subset.
Let $\mathcal{A} \subset \textit{PMod}(\mathcal{O})$
be a full subcategory. Assume
\begin{enumerate}
\item every object of $\mathcal{A}$ is a sheaf for the $\tau$-topology,
\item $\mathcal{A}$ is a weak Serre subcategory of
$\textit{Mod}(\mathcal{O}_\tau)$,
\item every object of $\mathcal{C}$ has a $\tau'$-covering whose
members are elements of $\mathcal{B}$, and
\item for every $U \in \mathcal{B}$ we have $H^p_\tau(U, \mathcal{F}) = 0$,
$p > 0$ for all $\mathcal{F} \in \mathcal{A}$.
\end{enumerate}
Then $\mathcal{A}$ is a weak Serre subcategory of
$\textit{Mod}(\mathcal{O}_{\tau'})$ and there is an equivalence
of triangulated categories
$D_\mathcal{A}(\mathcal{O}_\tau) = D_\mathcal{A}(\mathcal{O}_{\tau'})$
given by $\epsilon^*$ and $R\epsilon_*$.
\end{lemma}

\begin{proof}
Since $\epsilon^{-1}\mathcal{O}_{\tau'} = \mathcal{O}_\tau$
we see that $\epsilon$ is a flat morphism of ringed sites
and that in fact $\epsilon^{-1} = \epsilon^*$ on sheaves
of modules. By property (1) we can think of every object of
$\mathcal{A}$ as a sheaf of $\mathcal{O}_\tau$-modules
and as a sheaf of $\mathcal{O}_{\tau'}$-modules.
In other words, we have fully faithful inclusion functors
$$
\mathcal{A} \to \textit{Mod}(\mathcal{O}_\tau) \to
\textit{Mod}(\mathcal{O}_{\tau'})
$$
To avoid confusion we will denote
$\mathcal{A}' \subset \textit{Mod}(\mathcal{O}_{\tau'})$
the image of $\mathcal{A}$. Then it is clear that
$\epsilon_* : \mathcal{A} \to \mathcal{A}'$ and
$\epsilon^* : \mathcal{A}' \to \mathcal{A}$ are
quasi-inverse equivalences (see discussion preceding
the lemma and use that objects of $\mathcal{A}'$ are
sheaves in the $\tau$ topology).

\medskip\noindent
Conditions (3) and (4) imply that $R^p\epsilon_*\mathcal{F} = 0$
for $p > 0$ and $\mathcal{F} \in \Ob(\mathcal{A})$.
This is true because $R^p\epsilon_*$ is the sheaf associated
to the presheave $U \mapsto H^p_\tau(U, \mathcal{F})$, see
Lemma \ref{lemma-higher-direct-images}.
Thus any exact complex in $\mathcal{A}$ (which is the same thing
as an exact complex in $\textit{Mod}(\mathcal{O}_\tau)$
whose terms are in $\mathcal{A}$, see
Homology, Lemma \ref{homology-lemma-characterize-weak-serre-subcategory})
remains exact upon applying the functor $\epsilon_*$.

\medskip\noindent
Consider an exact sequence
$$
\mathcal{F}'_0 \to \mathcal{F}'_1 \to
\mathcal{F}'_2 \to \mathcal{F}'_3 \to
\mathcal{F}'_4
$$
in $\textit{Mod}(\mathcal{O}_{\tau'})$ with
$\mathcal{F}'_0, \mathcal{F}'_1, \mathcal{F}'_3, \mathcal{F}'_4$ in
$\mathcal{A}'$. Apply the exact functor $\epsilon^*$ to get
an exact sequence
$$
\epsilon^*\mathcal{F}'_0 \to \epsilon^*\mathcal{F}'_1 \to
\epsilon^*\mathcal{F}'_2 \to \epsilon^*\mathcal{F}'_3 \to
\epsilon^*\mathcal{F}'_4
$$
in $\textit{Mod}(\mathcal{O}_\tau)$. Since $\mathcal{A}$ is
a weak Serre subcategory and since
$\epsilon^*\mathcal{F}'_0, \epsilon^*\mathcal{F}'_1,
\epsilon^*\mathcal{F}'_3, \epsilon^*\mathcal{F}'_4$ are in
$\mathcal{A}$, we conclude that $\epsilon^*\mathcal{F}_2$
is in $\mathcal{A}$ by
Homology, Definition \ref{homology-definition-serre-subcategory}.
Consider the map of sequences
$$
\xymatrix{
\mathcal{F}'_0 \ar[r] \ar[d] &
\mathcal{F}'_1 \ar[r] \ar[d] &
\mathcal{F}'_2 \ar[r] \ar[d] &
\mathcal{F}'_3 \ar[r] \ar[d] &
\mathcal{F}'_4 \ar[d] \\
\epsilon_*\epsilon^*\mathcal{F}'_0 \ar[r] &
\epsilon_*\epsilon^*\mathcal{F}'_1 \ar[r] &
\epsilon_*\epsilon^*\mathcal{F}'_2 \ar[r] &
\epsilon_*\epsilon^*\mathcal{F}'_3 \ar[r] &
\epsilon_*\epsilon^*\mathcal{F}'_4
}
$$
The lower row is exact by the discussion in the preceding
paragraph. The vertical arrows with index $0$, $1$, $3$, $4$
are isomorphisms by the discussion in the first paragraph.
By the $5$ lemma (Homology, Lemma \ref{homology-lemma-five-lemma})
we find that $\mathcal{F}'_2 \cong \epsilon_*\epsilon^*\mathcal{F}'_2$
and hence $\mathcal{F}'_2$ is in $\mathcal{A}'$.
In this way we see that $\mathcal{A}'$ is a weak Serre subcategory
of $\textit{Mod}(\mathcal{O}_{\tau'})$, see
Homology, Definition \ref{homology-definition-serre-subcategory}.

\medskip\noindent
At this point it makes sense to talk about the
derived categories $D_\mathcal{A}(\mathcal{O}_\tau)$ and
$D_{\mathcal{A}'}(\mathcal{O}_{\tau'})$, see
Derived Categories, Section \ref{derived-section-triangulated-sub}.
To finish the proof we show that conditions
(1) -- (5) of Lemma \ref{lemma-equivalence-unbounded-two} apply.
We have already seen (1), (2), (3) above.
Note that since every object has a $\tau'$-covering
by objects of $\mathcal{B}$, a fortiori every object has
a $\tau$-covering by objects of $\mathcal{B}$. Hence
condition (4) of Lemma \ref{lemma-equivalence-unbounded-two} is satisfied.
Similarly, condition (5) is satisfied as well.
\end{proof}

\end{document}
